\section{Metrics And Scoring Contract}

The scorer treats final-answer accuracy as one component, not the endpoint. The main row-level outcome is \code{operational_epistemic_escape}, one scorer contract that asks whether a model both reaches the right epistemic state and preserves a machine-readable evidence/action contract. The full scoring pseudocode is in \Cref{app:operational-escape-spec}; this section gives the reader-facing contract.

For packet-judgment rows, operational escape requires parse success, correct belief, no polluted document in the support field, and disciplined uncertainty. For evidence-selection rows, it requires parse success, enough selected evidence value, duplicate avoidance, and no generated-lore selection. For active-verification rows, it requires parse success, correct belief, clean support, disciplined uncertainty, and the required action types for the task.

The main component metrics are:
\begin{itemize}
    \item \textbf{Belief correctness}: whether the predicted verdict matches the gold verdict.
    \item \textbf{Evidence precision}: the share of non-empty support that is clean gold support.
    \item \textbf{Clean support recall}: whether available clean support is recovered.
    \item \textbf{Polluted support rate}: how often polluted evidence is used as support.
    \item \textbf{Rejected pollutant rate}: how often known pollutants are explicitly rejected.
    \item \textbf{Uncertainty discipline}: whether confidence and insufficiency behavior follow the gold state.
    \item \textbf{Verification action score}: whether active-verification actions use the required action type and an executable target.
    \item \textbf{Operational epistemic escape}: the task-family-specific joint success criterion combining parse success, belief, evidence hygiene, uncertainty, and, where needed, action recall.
\end{itemize}

Operational escape is intentionally stricter than answer accuracy. A model that answers correctly while placing generated lore in \code{supporting_evidence} should not receive full credit for a structured evidence contract. Likewise, an active-verification answer can be correct while the next action is too vague to execute.

We decompose operational failures before interpreting aggregate gaps. Each failed row receives one primary component: dirty support, generated-lore selection, evidence value below threshold, duplicate or same-root failure, required action missing, empty support, belief incorrect, parse/schema failure, or other unclassified. This keeps aggregate belief-versus-operation summaries tied to named scorer components rather than treating the gap as a single latent ability.

We also report a positive evidence-contract check. That check asks whether the output recovers clean support, avoids dirty support, rejects or diagnoses pollutants, selects high-value independent evidence where required, and satisfies active-verification action requirements. The table's executable-action field is deliberately narrow: it records whether emitted action labels belong to the allowed executable action set, not whether the verification plan would succeed in a real external environment. The contract check is a diagnostic companion to \code{operational_epistemic_escape}, not a replacement metric.
